\documentclass[10pt,a4paper]{article}
\usepackage[top=1.4cm, bottom=1.4cm, left=1.6cm, right=1.6cm, headheight=14pt, footskip=18pt]{geometry}
\usepackage[utf8]{inputenc}
\usepackage{amsmath,amsfonts,amssymb}
\usepackage{graphicx}
\usepackage{tikz}
\usepackage{circuitikz}
\usetikzlibrary{calc,patterns}
\usepackage[most]{tcolorbox}
\usepackage{enumitem}
\usepackage{fancyhdr}
\usepackage{booktabs}
\usepackage{tabularx}
\usepackage{colortbl}
\usepackage{hyperref}

% --- Palet Warna Resmi UNIB ---
\definecolor{unibblue}{RGB}{0, 32, 96}
\definecolor{unibgold}{RGB}{197, 160, 89}
\definecolor{unibgreen}{RGB}{34, 112, 60}
\definecolor{boxbg}{RGB}{248, 250, 253}
\definecolor{linegray}{RGB}{210, 215, 225}

% --- Pengaturan Header & Footer ---
\pagestyle{fancy}
\fancyhf{}
\lhead{\footnotesize\textbf{\color{unibblue}Universitas Bengkulu} \textbar\ Program Studi S1 Teknik Elektro}
\rhead{\footnotesize\color{gray}Persamaan Diferensial (Genap 2026)}
\lfoot{\footnotesize\color{gray}Problem Set OBE \textbar\ \textbf{Video}: \url{https://youtu.be/iIYM-WRaTiU} \textbar\ \href{https://www.youtube.com/playlist?list=PLS5oOZWXZeTw}{\color{unibblue}\textbf{Playlist}}}
\cfoot{\footnotesize\color{gray}Hal. \thepage\ dari 2}
\rfoot{\footnotesize\color{unibblue}\textbf{Minggu 6: Transformasi Laplace \& Sinyal Tr...}}
\renewcommand{\headrulewidth}{0.6pt}
\renewcommand{\footrulewidth}{0.4pt}

\begin{document}

% ==================== HALAMAN 1 ====================
\noindent\begin{minipage}[c]{0.68\textwidth}%
  {\large\textbf{\color{unibblue}PROBLEM SET (TUGAS MANDIRI TERSTRUKTUR)}}\\[1mm]
  {\normalsize\textbf{Mata Kuliah: Persamaan Diferensial (2 SKS)}}\\[0.5mm]
  {\small\textbf{Topik Minggu 6:} Transformasi Laplace \& Sinyal Transien Impuls Surja Petir}%
\end{minipage}%
\hfill%
\begin{minipage}[c]{0.30\textwidth}%
  \begin{tcolorbox}[colback=unibblue!5,colframe=unibblue,boxrule=0.6pt,arc=2pt,left=4pt,right=4pt,top=2pt,bottom=2pt]
    \scriptsize
    \textbf{Nama:} \dotfill\\[1.2mm]
    \textbf{NPM:} \dotfill\\[1.2mm]
    \textbf{Kelas/Klp:} \dotfill\\[1.2mm]
    \textbf{Nilai Final:} \hfill \textbf{/ 100}
  \end{tcolorbox}%
\end{minipage}\par

\vspace{1.5mm}

% Kotak Sub-CPMK & Pedoman Tugas Terstruktur
\begin{tcolorbox}[colback=boxbg,colframe=unibgold!90!black,boxrule=0.6pt,arc=2pt,left=5pt,right=5pt,top=3pt,bottom=3pt,title=\textbf{\scriptsize Capaian Pembelajaran (Sub-CPMK 4) \& Pedoman Pengerjaan Tugas}]
  \scriptsize
  \textbf{Sub-CPMK 4:} Mahasiswa mampu memetakan fungsi waktu ke domain frekuensi kompleks $s$ menggunakan Transformasi Laplace serta memodelkan respons sistem terhadap sinyal diskontinu (Heaviside \& Dirac).\\
  \textbf{Video Perkuliahan (YouTube):} \url{https://youtu.be/iIYM-WRaTiU} \hfill \textbf{Playlist:} \href{https://www.youtube.com/playlist?list=PLS5oOZWXZeTw}{\color{unibblue}\textbf{bit.ly/pd-unib-playlist}}\\\textbf{Ketentuan Tugas:} Kerjakan secara mandiri pada kertas folio/A4 bergaris atau laporan digital rapi. Lampirkan penurunan rumus analitis lengkap dan cetak visualisasi kode program Julia (Soal C6). Kumpulkan pada awal perkuliahan minggu berikutnya.
\end{tcolorbox}

\vspace{1.5mm}

\noindent{\textbf{\color{unibblue}\large BAGIAN I: FONDASI KONSEPTUAL \& PENURUNAN MATEMATIS}}\par
\vspace{1mm}

% Soal C1
\noindent\textbf{1. (C1 - Mengingat \textbar\ Bobot: 10 Poin)}\par\vspace{0.8mm}
{\footnotesize Tuliskan definisi integral dari Transformasi Laplace unilateral $F(s) = \int_0^\infty f(t) e^{-st} dt$, dan tuliskan pasangan transformasi untuk tiga fungsi dasar: fungsi tangga satuan $u(t)$, fungsi impuls Dirac $\delta(t)$, dan fungsi eksponensial $e^{-at}$.}\par

\vspace{2.5mm}

% Soal C2
\noindent\textbf{2. (C2 - Memahami \textbar\ Bobot: 15 Poin)}\par\vspace{0.8mm}
{\footnotesize Jelaskan Teorema Pergeseran Frekuensi kompleks (*First Shifting Theorem*): $\mathcal{L}\{e^{-at} f(t)\} = F(s+a)$. Bagaimana sifat pergeseran frekuensi ini menjelaskan redaman eksponensial pada respons transien rangkaian osilator teredam?}\par

\vspace{2.5mm}

% Soal C3
\noindent\textbf{3. (C3 - Menerapkan \textbar\ Bobot: 20 Poin)}\par\vspace{0.8mm}
\noindent\begin{minipage}[t]{0.60\textwidth}%
  {\footnotesize Berdasarkan standar pengujian tegangan tinggi \textbf{IEC 60060-1}, gelombang surja impuls petir standar $1{,}2/50\,\mu\text{s}$ dinyatakan sebagai fungsi eksponensial ganda: $v(t) = V_0 (e^{-\alpha t} - e^{-\beta t}) u(t)$. Tentukan representasi aljabar domain Laplace $V(s) = \mathcal{L}\{v(t)\}$ secara eksplisit.}%
\end{minipage}%
\hfill%
\begin{minipage}[t]{0.37\textwidth}%
  \centering
  \resizebox{0.95\linewidth}{!}{%
  \begin{circuitikz}[american, scale=0.68, transform shape]
    \draw (0,0) to[I, l=$I_0 \delta(t)$] (0,1.8)
          to[short] (1.5,1.8)
          to[R, l=$R_1$] (1.5,0)
          to[short] (0,0);
    \draw (1.5,1.8) to[R, l=$R_2$] (3.0,1.8)
          to[C, l=$C_2$, v=$v_{\text{impuls}}(t)$] (3.0,0)
          to[short] (1.5,0);
  \end{circuitikz}%
  }%
\end{minipage}\par

\vfill

% Kotak Formula Rujukan Teknis & Standar Industri
\begin{tcolorbox}[colback=unibblue!3,colframe=unibblue,colbacktitle=unibblue,coltitle=white,boxrule=0.6pt,arc=2pt,left=6pt,right=6pt,top=3pt,bottom=3pt,title=\textbf{\scriptsize FORMULA RUJUKAN TEKNIS \& STANDAR INDUSTRI TERKAIT}]
  \scriptsize
  \begin{itemize}[leftmargin=*, nosep]
  \item Pasangan Transformasi Dasar: $\mathcal{L}\{u(t)\} = \frac{1}{s}$, $\mathcal{L}\{\delta(t)\} = 1$, $\mathcal{L}\{e^{-at}\} = \frac{1}{s+a}$, $\mathcal{L}\{\sin\omega t\} = \frac{\omega}{s^2+\omega^2}$
  \item Teorema Translasi Frekuensi \& Waktu: $\mathcal{L}\{e^{-at} f(t)\} = F(s+a)$ \quad\textbar\quad $\mathcal{L}\{f(t-T)u(t-T)\} = e^{-sT} F(s)$
  \item Teorema Nilai Awal \& Akhir: $\lim_{t\to 0^+} f(t) = \lim_{s\to\infty} s F(s)$ \quad\textbar\quad $\lim_{t\to\infty} f(t) = \lim_{s\to 0} s F(s)$
  \item Standar \textbf{IEC 60060-1}: Impuls Petir Standar $1{,}2/50\,\mu\text{s}$, $v(t) = V_0(e^{-\alpha t} - e^{-\beta t})$
\end{itemize}
\end{tcolorbox}

\newpage
% ==================== HALAMAN 2 ====================

\noindent{\textbf{\color{unibblue}\large BAGIAN II: ANALISIS SISTEM, EVALUASI STANDAR, \& KOMPUTASI JULIA}}\par
\vspace{1mm}

% Soal C4
\noindent\textbf{4. (C4 - Menganalisis \textbar\ Bobot: 20 Poin)}\par\vspace{0.8mm}
{\footnotesize Analisis respons transien rangkaian RL terhadap tegangan pulsa persegi berdurasi terbatas $T$: $v_s(t) = V_0 [u(t) - u(t-T)]$. Gunakan Teorema Pergeseran Waktu (*Second Shifting Theorem*) untuk menurunkan arus domain-$s$ $I(s)$, lalu rekonstruksi bentuk gelombang arus $i(t)$ untuk interval $0 \le t < T$ dan interval $t \ge T$.}\par

\vspace{3.5mm}

% Soal C5
\noindent\textbf{5. (C5 - Mengevaluasi \textbar\ Bobot: 20 Poin)}\par\vspace{0.8mm}
{\footnotesize Evaluasi perilaku sistem domain frekuensi menggunakan Teorema Nilai Awal (IVT) dan Teorema Nilai Akhir (FVT): Diberikan fungsi transfer tegangan kapasitor $V_C(s) = \frac{100}{s(s^2 + 6s + 25)}$. Tentukan nilai awal $v_C(0^+)$ dan nilai akhir tunak $v_C(\infty)$ secara langsung dari domain-$s$, lalu validasi hasilnya dengan penyelesaian domain waktu.}\par

\vspace{3.5mm}

% Soal C6
\noindent\textbf{6. (C6 - Merancang \& Komputasi Julia \textbar\ Bobot: 15 Poin)}\par\vspace{0.8mm}
{\footnotesize Rancang skrip program ilmiah berbasis \textbf{Julia} untuk menghasilkan dan memplot gelombang impuls petir standar IEC 60060-1 ($V_0 = 250\text{ kV}, \alpha = 14658\text{ s}^{-1}, \beta = 2469135\text{ s}^{-1}$) pada selang waktu mikrodetik $t \in [0, 80\,\mu\text{s}]$. Tandai titik waktu muka (*front time* $t_1 = 1{,}2\,\mu\text{s}$) dan waktu paruh (*tail time* $t_2 = 50\,\mu\text{s}$).}\par

\vfill

% Tabel Rekapitulasi Skor OBE & Lembar Catatan Umpan Balik Dosen
\noindent\begin{minipage}[b]{0.62\textwidth}%
  \centering
  \scriptsize
  \begin{tabular}{|c|c|c|c|c|c|c|}
    \hline
    \rowcolor{unibblue}
    \textbf{\color{white}C1} & \textbf{\color{white}C2} & \textbf{\color{white}C3} & \textbf{\color{white}C4} & \textbf{\color{white}C5} & \textbf{\color{white}C6} & \textbf{\color{white}Total} \\
    \rowcolor{unibblue!10}
    10 Pts & 15 Pts & 20 Pts & 20 Pts & 20 Pts & 15 Pts & 100 Pts \\
    \hline
    & & & & & & \\[3mm]
    \hline
  \end{tabular}%
\end{minipage}%
\hfill%
\begin{minipage}[b]{0.35\textwidth}%
  \centering
  \scriptsize
  \textbf{Verifikasi Dosen / Asisten:}\\[6mm]
  (\dotfill)\\[1mm]
  \tiny Tanggal: \dotfill
\end{minipage}\par

\vspace{2mm}
\begin{tcolorbox}[colback=white,colframe=unibblue!40!black,colbacktitle=unibblue!10,coltitle=unibblue,boxrule=0.5pt,arc=2pt,left=5pt,right=5pt,top=3pt,bottom=3pt,title=\textbf{\tiny Catatan Evaluasi \& Umpan Balik Dosen Pengampu}]
  \tiny
  \textbf{Rubrik Pemeriksaan Pengerjaan Tugas Mandiri:}\par\vspace{0.8mm}
  $\square$ Penurunan matematis langkah-demi-langkah runtut (C1--C4) \hfill $\square$ Validasi analitis \& standar industri IEEE/IEC/SPLN akurat (C5)\\[0.8mm]
  $\square$ Skrip program Julia mandiri \& grafik visualisasi terlampir (C6) \hfill $\square$ Kerapian format laporan \& ketepatan waktu pengumpulan\\[2.0mm]
  \textbf{Catatan / Koreksi Khusus Dosen:}\\[1.6cm]
\end{tcolorbox}

\end{document}
